\chapter{Ce qu'on a mesuré sur le lab}
\label{ch:resultats}

Ce chapitre présente les chiffres obtenus en faisant tourner
HOMO-CI sur le lab UrbanUpC. L'objectif est de répondre aux trois
questions posées en introduction~: couverture, fraîcheur,
précision. Les résultats sont donnés avec leurs limites~; il n'y
a pas de p-value, juste des mesures avec leur protocole et leur
sensibilité aux choix faits.

\section{Protocole en bref}

\paragraph{Cible.} Le lab UrbanUpC complet~: les 3 VMs Azure, le
domaine \texttt{corpnet.local}, les 3 applications, le SIEM
Wazuh.

\paragraph{Référence pour la couverture.} Un catalogue de 100
preuves attendues, construit par lecture du \emph{Guide
d'homologation ANSSI}~\cite{anssi-homologation-2014} et de la
norme \emph{ISO/IEC 27002:2022}. Chaque preuve est classée par
étape ANSSI et marquée comme automatisable (A), semi-automatisable
(B), ou intrinsèquement manuelle (C).

\paragraph{Période de mesure.} Trois semaines d'exécutions
nocturnes (\emph{nightly}) du pipeline, du 22~avril au 13~mai
2026, sur le lab. Plus deux injections \emph{ad hoc} pour
mesurer la latence (création et suppression d'une règle NSG
chronométrées).

\section{Couverture des preuves}

\paragraph{Question.} Sur les 100 preuves attendues, combien
HOMO-CI arrive-t-il à collecter ?

\paragraph{Résultat brut.} L'outil collecte automatiquement
\textbf{40 preuves} sur 100. Si on pondère par
\emph{automatisabilité} (A~=~1,0~; B~=~0,5~; C~=~0,0), le score
est de \textbf{49\,\%}.

C'est en dessous du seuil de 60\,\% qu'on s'était fixé. La raison
saute aux yeux quand on regarde la décomposition par étape
ANSSI~:

\begin{table}[h]
\centering
\caption{Couverture par scope d'étapes ANSSI.}
\label{tab:couverture}
\small
\begin{tabular}{lrrl}
\toprule
\textbf{Scope} & \textbf{Items} & \textbf{Couverture pondérée} & \textbf{Verdict} \\
\midrule
Technique (étapes 4, 6, 7) & 62 & \textbf{61,5\,\%} & Atteint l'objectif \\
Gouvernance (étapes 1, 2, 3, 5, 8) & 26 & 11,5\,\% & Non automatisable \\
Maintien (étape 9) & 12 & 45,8\,\% & Partiel \\
\bottomrule
\end{tabular}
\end{table}

\paragraph{Lecture.} Sur les preuves \emph{techniques} (règles
NSG, CVE, comptes AD, règles SIEM), HOMO-CI dépasse 60\,\%. Sur
les preuves \emph{de gouvernance} (formation, sensibilisation,
revues humaines), il ne fait presque rien --- et c'est attendu~:
ces preuves sont par nature manuelles. L'outil sert à couvrir
la moitié technique du dossier, pas à se substituer aux
décisions humaines.

\paragraph{Sensibilité.} On a refait la classification A/B/C
avec deux variantes pessimiste (faire glisser 10 items de B vers
C) et optimiste (10 items de B vers A)~:

\begin{table}[h]
\centering
\caption{Sensibilité à la classification A/B/C (scope technique).}
\label{tab:sensibilite-cov}
\small
\begin{tabular}{lr}
\toprule
\textbf{Variante} & \textbf{Couverture pondérée} \\
\midrule
Référence & 61,5\,\% \\
Pessimiste & 53,2\,\% \\
Optimiste & 67,7\,\% \\
\bottomrule
\end{tabular}
\end{table}

Le résultat est robuste mais sensible à la rigueur de
classification~: une vraie validation passerait par une
classification croisée à deux personnes.

\section{Fraîcheur : combien de temps pour propager un changement}

\paragraph{Question.} Quand une règle NSG change réellement,
combien de temps faut-il pour que le dossier d'homologation
généré par HOMO-CI reflète ce changement~?

\paragraph{Modèle théorique.} Si le pipeline tourne en nightly,
on s'attend à un délai moyen de 12 heures (au pire 24~h, dans le
cas où le changement arrive juste après une exécution).

\paragraph{Mesure réelle sur Azure.}
Deux injections chronométrées~:

\begin{table}[h]
\centering
\caption{Latence de propagation mesurée (deux injections sur
\texttt{nsg-dmz}).}
\label{tab:latence}
\small
\begin{tabular}{lrr}
\toprule
\textbf{Mesure} & \textbf{Détection (s)} & \textbf{Dossier régénéré (s)} \\
\midrule
Création d'une règle & 2,99 & 4,99 \\
Suppression d'une règle & 5,04 & 7,05 \\
\midrule
\textbf{Moyenne} & \textbf{4,0\,s} & \textbf{6,0\,s} \\
\bottomrule
\end{tabular}
\end{table}

Évidemment, ces chiffres correspondent à un mode \emph{déclenché
manuellement}~: l'opérateur lance \texttt{homo-ci collect} juste
après le changement. En mode nightly réel, on retombe sur les
12--24~heures attendues.

\paragraph{Comparaison.} Un cycle d'audit annuel place ce délai à
plusieurs mois. Même en mode nightly, HOMO-CI est largement plus
rapide. Le tableau ci-dessous montre l'effet du choix de
cadence~:

\begin{table}[h]
\centering
\caption{Latence en fonction de la cadence de régénération.}
\label{tab:cadence}
\small
\begin{tabular}{lrr}
\toprule
\textbf{Cadence} & \textbf{Délai moyen} & \textbf{P95} \\
\midrule
Weekly & 3,8\,j & 6,7\,j \\
Nightly (référence) & 13,5\,h & 22,8\,h \\
Toutes les 6\,h & 3,7\,h & 5,9\,h \\
Toutes les heures & 38\,min & 58\,min \\
Toutes les 15\,min & 8\,min & 14\,min \\
\bottomrule
\end{tabular}
\end{table}

Conclusion pratique~: la cadence est un levier simple. Pour un
ENT en production, du nightly est probablement suffisant. Pour
une infra plus critique, on peut descendre à l'heure sans
modifier le code, juste en changeant le \texttt{cron}.

\section{Précision : ce que dit le dossier correspond-il à la réalité}

\paragraph{Question.} Un dossier généré aujourd'hui décrit-il
fidèlement le lab tel qu'il est ce matin~?

\paragraph{Protocole.} Simulation rétroactive sur 180 jours avec
125 éléments de SI (NSG, CVE, comptes AD, règles SIEM) qui
évoluent selon un modèle de changement réaliste (basé sur le
journal git du lab). On compare le \emph{taux d'alignement}
entre~:

\begin{itemize}
  \item un dossier statique généré à $t_0$ et jamais mis à jour~;
  \item un dossier régénéré chaque nuit par HOMO-CI.
\end{itemize}

\paragraph{Résultats.}

\begin{table}[h]
\centering
\caption{Taux d'alignement par âge du dossier (simulation 180\,j).}
\label{tab:alignement}
\small
\begin{tabular}{lrr}
\toprule
\textbf{Dossier} & \textbf{Âge (jours)} & \textbf{Alignement} \\
\midrule
Statique $D_0$ & 0 & 100\,\% \\
Statique $D_{30}$ & 30 & 83,3\,\% \\
Statique $D_{90}$ & 90 & \textbf{63,1\,\%} \\
Statique $D_{150}$ & 150 & 51,1\,\% \\
Statique $D_{180}$ & 180 & 46,5\,\% \\
\midrule
HOMO-CI (nightly) & --- & \textbf{100\,\%} \\
\bottomrule
\end{tabular}
\end{table}

\paragraph{Lecture.} À 90 jours, un dossier statique perd
36\,points d'alignement~: plus d'un tiers de ce qu'il décrit ne
correspond plus à la réalité. HOMO-CI maintient l'alignement à
100\,\%, par construction (il ne décrit que ce qu'il a collecté
la nuit précédente).

\todofig{Courbe d'alignement en fonction de l'âge du dossier
(statique en rouge décroissant, continu en bleu constant à
100\,\%).}

\paragraph{Honnêteté sur le 100\,\%.} Ce 100\,\% ne signifie pas
\emph{parfait}. Cela signifie~: \emph{ce que le dossier décrit
est ce que l'outil a vu cette nuit}. Si l'outil a un trou de
couverture (cas des preuves de gouvernance vu plus haut), ce trou
reste un trou. On parle d'alignement \emph{sur le périmètre que
l'outil couvre}.

\section{Coût de l'outil}

\paragraph{Coût Azure.} Sur les trois semaines de test, le crédit
consommé par les VMs et le Blob Storage est resté sous
275\,\euro. Le pipeline lui-même tourne sur le poste local et ne
consomme pas de ressources cloud.

\paragraph{Temps d'exécution.} Un cycle complet (\texttt{collect},
\texttt{diff}, \texttt{render}, \texttt{publish}) prend environ
90~secondes sur le lab,
dont 70\,\% sur les appels API Azure. Pour un SI plus gros,
l'élargissement viendrait surtout du nombre de NSG et de l'AD.

\paragraph{Temps de mise en place.} Cloner le dépôt et faire
tourner une première exécution sur un nouveau lab prend environ
4~heures (création du Resource Group, déploiement des VMs,
installation de Wazuh, configuration des credentials Azure).

\section{Synthèse}

\begin{table}[h]
\centering
\caption{Synthèse des trois objectifs.}
\label{tab:synthese}
\small
\begin{tabularx}{\textwidth}{lXX}
\toprule
\textbf{Objectif} & \textbf{Résultat} & \textbf{Verdict} \\
\midrule
Couverture des preuves & 61,5\,\% (scope technique), 49\,\% (global) & Atteint sur le périmètre technique, comme prévu \\
Fraîcheur du dossier & 4--6\,s en mode déclenché, 13,5\,h en nightly & Très en-dessous des 24\,h visées \\
Précision du dossier & 100\,\% sur le périmètre couvert, vs 63\,\% pour un dossier statique de 90\,j & Net gain sur la dérive \\
\bottomrule
\end{tabularx}
\end{table}

\begin{takeaway}
L'outil fait ce qu'on attendait sur la couche technique. Sur la
gouvernance, il est inutile, et c'est volontaire. Le gain
principal n'est pas dans l'innovation algorithmique, mais dans
la régularité~: tous les jours, à 03h00, un nouveau dossier est
produit, daté, signé, comparable au précédent.
\end{takeaway}
